\documentclass{article}
\usepackage{amsmath}
\usepackage{graphicx}
\begin{document}
\section{Math}
Mapping a dataflow graph onto a coarse-grained array asks the compiler to place every operation and route every value in space and time, and the search grows with the fabric.
\begin{align}
 a &= b + c \label{e:a}\\
 d &= e + f \label{e:d}\\
 d' &= e \label{e:dd}\\
 g &= \text{the new bound} \label{e:g}
\end{align}
See \eqref{e:a}, \eqref{e:d}, \eqref{e:h} and \eqref{e:k}.
\begin{gather*}
 x = 1 \\
 y = 3
\end{gather*}
Mapping a dataflow graph onto a coarse-grained array asks the compiler to place every operation and route every value in space and time, and the search grows with the fabric.
Mapping a dataflow graph onto a coarse-grained array asks the compiler to place every operation and route every value in space and time, and the search grows with the fabric.\footnote{A note with $\alpha - \beta$ in it.}
\begin{equation}\label{e:k}
 k = 0
\end{equation}
\begin{figure}[h]\centering\includegraphics[width=.2\linewidth]{example-image-a}\caption{Cost $c = n \log n$ per step.}\end{figure}
Mapping a dataflow graph onto a coarse-grained array asks the compiler to place every operation and route every value in space and time, and the search grows with the fabric.
\end{document}
